\documentclass[10pt,a4paper]{amsart}
\usepackage[margin=19mm]{geometry}
\usepackage{amsmath,amssymb,mathtools}
\usepackage[scr=rsfs,cal=euler]{mathalfa}
\usepackage{xfrac}
\usepackage[unicode,hidelinks,bookmarks=false]{hyperref}
\newcommand{\ie}{i.e.\ }
\newcommand{\cf}{cf.\ }
\newcommand{\resp}{respectively\ }
\newcommand{\Subsection}{\subsection}
\providecommand{\nobreakdash}{\nobreak\hbox{-}}
\setlength{\emergencystretch}{2em}
\multlinegap0pt
\usepackage{bm}
\usepackage{bbm}
\usepackage{dsfont}
\usepackage{xspace}
\usepackage{cancel}
\usepackage{mathtools}
\usepackage{amsthm}
\makeatletter
\@ifundefined{shared}{\newtheorem{shared}{Shared}}{}
\@ifundefined{lemma}{\newtheorem{lemma}[shared]{Lemma}}{}
\makeatother
\title[Critical percolation on preferential attachment graphs with ]{Critical percolation on preferential attachment graphs with infinite variance}
\author{Peter M\"orters, Lucas Sch\"atze}
\date{Source: arXiv:2606.27844v1 (CC BY 4.0)}
\begin{document}
\maketitle

\noindent\emph{Source and licence.} Critical percolation on preferential attachment graphs with infinite variance. Version arXiv:2606.27844v1 (CC BY 4.0), distributed under the Creative Commons Attribution 4.0 International licence, as recorded in the arXiv licence metadata for this exact version and shown on the abstract page https://arxiv.org/abs/2606.27844v1. Full rights notice: licenses/arxiv-2606.27844-CC-BY-4.0.txt.\par\smallskip

\noindent\textit{Editorial scope.} Counted unit: the Lemma whose source label is \texttt{pathlength} in the article (source0.tex lines 673--683, source label \texttt{pathlength}), delivered together with the article's own complete original proof of it (source0.tex lines 684--788, one \texttt{proof} environment of 105 lines). No mathematics is written, repaired or re-derived: statement and proof are the article's own text line for line. Required context retained uncounted: mmudef (lines 467--469); 1.3 (lines 660--671). Every definition, display and label of those lines is retained unchanged. Omitted from this package and not redistributed: 1--466, 470--659, 672--672, 789--1725. Nothing inside a retained excerpt is omitted or altered: every retained body is delivered exactly as the article's lines are written, so no omission receipt of any kind accompanies this package. Citations: the retained excerpts carry no citation command at all, so no bibliography entry of the article is transcribed and no reference is redistributed. Reference adaptation: none was needed and none was made; every reference inside the delivered unit resolves inside it, so no unresolved reference or citation is produced. Printed slips kept literal: no printed word, symbol, hypothesis or number of the counted statement or of its proof was altered. Layout and class: the article's own class is \texttt{article} with packages including inputenc, fontenc, babel, wrapfig, amsmath, bbm, xcolor, centernot, amsfonts, enumitem, amsthm, mathrsfs, amssymb, hyperref; the wrapper is the lane's pinned \texttt{amsart} editorial wrapper at 10pt on A4, which restates the theorem-like environments the retained text needs and adds the article's own macro definitions that the retained text needs (none), with extra wrapper packages (bm, bbm, dsfont, xspace, cancel, mathtools). The delivered numbering is the unit's own. Assets: no retained span contains a figure environment, a graphics-inclusion command, a PDF-page inclusion, a table or a TikZ picture; no image file is copied into this package, the package declares no asset dependency and no image file is redistributed with it. Licence: version arXiv:2606.27844v1 of this article is distributed under the Creative Commons Attribution 4.0 International licence, as recorded in the arXiv licence metadata for this exact version and shown on the abstract page https://arxiv.org/abs/2606.27844v1. A full rights notice is delivered at licenses/arxiv-2606.27844-CC-BY-4.0.txt. Characters: every retained excerpt is pure ASCII, so the delivered engine, XeLaTeX with its default Latin Modern text font, reports no missing character.
\begin{align}\label{mmudef}
    m:=\lfloor \mu \,n\,\beta(n)^{\alpha}\rfloor & \text{ and } s:=\beta(n)^{\alpha /2 } \text{  for } \alpha:=\tfrac{2}{2\gamma-1},
\end{align} 

    \begin{lemma} \label{1.3}
		Let $\mu \colon\{1,\dots,n\}\to[0,\infty)$ be a function satisfying, for some $\psi, \phi >0$ and $s_1\in(0,1)$, that
		\begin{align*}
			\mu(x)\le \mathbbm{1}_{\{x\ge s_1 n\}}\psi x^{\gamma-1}+\phi x^{-\gamma}
            \qquad
		\text{for all $x\in\{1,\dots,n\}$. }
        \end{align*}
        Then, for any $s_2\le s_1 $,  we have that
		\begin{align*}
			\sum_{s_2 n<y\le n}\mu(y) p_{xy} & \le \mathbbm{1}_{\{x >s_2 n\}}\beta(n) \Big(\phi \,\frac{(s_2 n)^{1-2\gamma}}{2 \gamma -1}+\psi\log\left(\tfrac{1}{s_2}\right)\Big) x^{\gamma-1}\\& \phantom{xxxx} +\beta(n) \Big(\phi \log\left(\tfrac{1}{s_2}\right)+\psi\,\frac{n^{2 \gamma -1}}{2 \gamma -1}\Big) x^{-\gamma}.
		\end{align*}
	\end{lemma}

\begin{lemma}\label{pathlength}
Let
$\mathcal{N}_k(\mathcal G_{n,\mu}^{_{(2)}})$ the set of vertices $x\in\{m+1,\ldots,n\}$, which are connected to $\mathcal G_{n,\mu}^{_{(2)}}$ by a shortest path of length at least $k$.  Then for any $k\ge\big\lfloor\tfrac{1-\gamma}{2\gamma-1}\big\rfloor +1  $ we have 
$$\frac{\#\mathcal{N}_{2k+1}(\mathcal G_{n,\mu}^{(2)})}{n\beta(n)^{\alpha/2}}\overset{\mathbb P}{\longrightarrow }0\qquad \text{ in probability.}$$
Also, all vertices at even distances are negligible, that is
$$\frac{1}{n\beta(n)^{\alpha/2}}\sum_{k=1}^\infty \#\big(\mathcal{N}_{2k}(\mathcal G_{n,\mu}^{(2)})\setminus \mathcal{N}_{2k+1}(\mathcal G_{n,\mu}^{(2)})\big)
\longrightarrow 0\qquad \text{ in probability.}$$
In particular if $\gamma>\tfrac23$, then
$$\frac{\#\mathcal{N}_2(\mathcal G_{n,\mu}^{(2)})}{n\beta(n)^{\alpha/2}}\longrightarrow 0\qquad \text{ in probability.}$$
    
\end{lemma}

\begin{proof}[Proof of Lemma~\ref{pathlength}]
The proof uses a path counting argument. Recall $m, s$ from \eqref{mmudef} %Let $m=\lfloor \mu \,n\,\beta(n)^{\alpha}\rfloor$, 
and let $\tilde s=\tilde s(n)$ to be determined, but such that $m \le \tilde sn \ll sn $.
Let $X_l$ be the number of paths $x_0x_1\cdots x_{l-1}x_l$ of length $l$, which satisfy the following conditions 
\begin{enumerate}
    \item[(a)] If $l=2$, then $x_0,x_1>m$ and $x_2 \le m$.
    \item[(b)] If $l\ge 3$, then $x_k>\tilde sn$ for all $k<l$ and $x_l\le \tilde sn$. 
\end{enumerate}
Then we have the following upper bound 
\begin{align} \label{13}
   \frac{1}{sn}\#\mathcal{N}_k(\mathcal G_{n,\mu}^{(2)}) \le    \frac{1}{sn}\sum_{l\ge k} X_l + o(1),
\end{align}
using that the total number of vertices in $\{m+1,\ldots, \tilde{s}n\}$ is $o(sn)$.
It remains to show that the other terms in (\ref{13}) go to zero in probability. We have
\begin{align*}
    \mathbb{E}[X_2]=&\sum_{i=1}^m \sum_{x,y=m+1}^n \beta(n)^2 i^{-\gamma}x^{\gamma-1}(x\wedge y)^{-\gamma}(x\vee y)^{\gamma-1}
    \\ \le & \sum_{i=1}^m \sum_{x,y=m+1}^n \beta(n)^2 i^{-\gamma}\big(x^{-1}y^{\gamma-1}+y^{-\gamma}x^{2(\gamma-1)}\big)
   \\ \le &  C\beta(n)^2 m^{1-\gamma}\big(\log(\tfrac{n}{m})n^{\gamma}+n^{\gamma}\big)
    \le   C n \beta(n)^{\gamma\alpha}\log(\beta(n)^{-\alpha}) = o(sn)\,,
\end{align*}
as $\gamma>\tfrac12$. 
For $l\ge 3$ we look at the expectation
\begin{align*}
	     \mathbb{E}[X_l]&\le \sum_{0<x_l\le \tilde sn} \; \sum_{\tilde sn<x_{l-1}\le n} \cdots \sum_{\tilde sn<x_{0}\le n} \prod_{i=0}^{l-1} \beta(n)(x_i \vee x_{i+1})^{\gamma-1}(x_i \wedge x_{i+1})^{-\gamma} \; .
	\end{align*}
	Note that in the above we would actually get an equality if we restrict the summation variables to be distinct.
	We define, for all $k\in\{1,\dots,l\}$,
	\begin{align*}
		\mu_k(x_k):=\sum_{\tilde sn<x_{k-1}\le n} \cdots \sum_{\tilde sn<x_{0}\le n} \prod_{i=0}^{k-1} \, \beta(n)(x_i \vee x_{i+1})^{\gamma-1}(x_i \wedge x_{i+1})^{-\gamma}.
	\end{align*}
	This is a recursion, since
\begin{align} \label{recurmu}
  \mu_{k+1}(x)=\sum_{\tilde sn<y\le n}\mu_k(y) \beta(n) (x \vee y)^{\gamma-1}(x \wedge y)^{-\gamma} \quad \text{ for all } k\ge1 \; .
\end{align}
Thus we can write for the expectation
\begin{align} \label{14}
		\mathbb{E}[X_l]&\le \sum_{0<x_l\le \tilde sn} \mu_l(x_l)
\end{align}
and apply Lemma \ref{1.3}. For the involved variables  we choose $s_1=s_2=\tilde s =\beta(n)^{r}$, with $r\in(\tfrac{1}{2\gamma-1},\tfrac{2}{2\gamma-1}]$. This guarantees $m\le \tilde s n \ll sn$.
 First we have for the starting value of our recursion
\begin{align*}
    \mu_1(x_1)= & \sum_{\tilde sn<x_{0}\le n} \beta(n)(x_0 \vee x_{1})^{\gamma-1}(x_0 \wedge x_{1})^{-\gamma}
		\\ \le & \mathbbm{1}_{\{x_1 >\tilde sn\}}\frac{\beta(n) n^{1-\gamma}}{1-\gamma}\, x_{1}^{\gamma-1}+\frac{\beta(n) n^{\gamma}}{\gamma}\, x_{1}^{-\gamma} \;.
\end{align*}
Thus we define $\psi_1=\tfrac{\beta(n) n^{1-\gamma}}{1-\gamma}$ and $\phi_1=\tfrac{\beta(n) n^{\gamma}}{\gamma}$. Then, using Lemma \ref{1.3}, we define recursively for all $k\ge1$,
\begin{equation}
    \label{psiphi}
\begin{aligned}
		\psi_{k+1}&=\phi_k  \beta(n)\frac{(\tilde s n)^{1-2\gamma}}{2\gamma-1}+\psi_k \beta(n) \log\left(\tfrac{1}{\tilde s}\right),
		\\ \phi_{k+1}&=\phi_k \beta(n) \log\left(\tfrac{1}{\tilde s}\right)+\psi_k \beta(n) \frac{n^{2 \gamma -1}}{2\gamma-1} \;.
\end{aligned}
\end{equation}
From this definition one can see immediately
\begin{align*}
    \psi_2 =O\Big(\beta(n)^{2}\tilde s ^{1-2\gamma} n^{1-\gamma}  \Big) \quad \text{ and } \quad
    \phi_2 =O\Big( \beta(n)^2  \log\left(\tfrac{1}{\tilde s}\right) n^{\gamma}  \Big) \;.
\end{align*}
However, one should calculate the first few terms. We give one more
\begin{align*}
    \psi_3 =O\Big(\beta(n)^{3}\tilde s ^{1-2\gamma}\log\left(\tfrac{1}{\tilde s}\right) n^{1-\gamma}  \Big) \quad \text{ and } \quad
    \phi_3 =O\Big( \beta(n)^{3}\tilde s ^{1-2\gamma}   n^{\gamma}  \Big)  \;.
\end{align*}
One can see by induction that, for all $k\ge1$,
\begin{equation}
\begin{aligned} \label{15}
   \psi_k =&\,O\Big(\beta(n)^{k}\tilde s ^{\lfloor \tfrac{k}{2}\rfloor(1-2\gamma)} \log\left(\tfrac{1}{\tilde s}\right)^{\mathbbm{1}_{\{\text{k odd}\}}} n^{1-\gamma} \Big), \\
    \phi_k =&\,O\Big(\beta(n)^{k}\tilde s ^{( \lceil \tfrac{k}{2}\rceil -1)(1-2\gamma)} \log\left(\tfrac{1}{\tilde s}\right)^{\mathbbm{1}_{\{\text{k even}\}}} n^{\gamma}  \Big) .
\end{aligned}    
\end{equation}
Observe that $\phi_k=o(\phi_{k-1})$ if $k$ is even. 
%Thus the sum over all path-lengths reduces to the sum over the paths of uneven length. {\color{red}Say explicitly what this means.}{\color{blue}this is exactly what happens below} Now b
By (\ref{14}) and Lemma \ref{1.3} we get
\begin{align*}
    \sum_{l=3}^n \mathbb{E}[X_l] &\le \sum_{l=1}^\infty 2\sum_{0<x_{2l+1}\le \tilde s n} \phi_{2l+1} x_{2l+1}^{-\gamma}
    \le c n  \tilde s^{1-\gamma}\beta(n)^3 \tilde s ^{1-2\gamma}\sum_{l=0}^\infty \big( \beta(n)^{2 }\tilde s ^{(1-2\gamma)}\big)^l\,.
\end{align*}
Using that $\tilde s = \beta(n)^r$ we see that the infinite sum above converges if and only if $r<\alpha$, as %should be 
expected. If $\gamma>\tfrac23$ we can choose $r>\tfrac{1}{2\gamma-1}$, such that the term above is~$o(sn)$. For the sake of simplicity let \smash{$r=\tfrac{4}{3}\tfrac{1}{2\gamma-1}$}, then 
\begin{align*}
     \sum_{l=3}^n \mathbb{E}[X_l] &\le  \,C n \beta(n)^{(2\gamma-\tfrac13)\tfrac{1}{2\gamma-1}} = o(sn)\,,
\end{align*}
since $2\gamma-\tfrac13>1$ if $\gamma>\tfrac23$. Now if $\gamma\in(\tfrac12,\tfrac23]$ we can still reduce the $k$-neighbourhood of $\mathcal G^{_{(2)}}_{n,\mu}$ to certain finite $k$. Fix $r=\tfrac{1+\varepsilon}{2\gamma-1}$ for some small $\varepsilon>0$ and $k\ge1$. Then, similar to before, we have 
\begin{align*}
     \sum_{l=2k+1}^n \mathbb{E}[X_l] &\le \sum_{l=k}^\infty 2\sum_{0<x_{2l+1}\le \tilde s n} \phi_{2l+1} x_{2l+1}^{-\gamma}
     \le C n  \beta(n)^{2k+1} \tilde s ^{k(1-2\gamma)+1-\gamma} \,,
\end{align*}
which is $o(sn)$ if %and only if 
%\begin{align*}
    $2k+1+r(1-\gamma)+k(1-2\gamma)r>\tfrac{1}{2\gamma-1}$,
%  \\
%    \iff \quad (1-\varepsilon)k>&\tfrac{(1-\varepsilon)(1-\gamma)}{2\gamma-1}  \iff \quad 
and hence for $k\ge \lfloor\frac{1-\gamma}{2\gamma-1}\rfloor +1$.% 
\smallskip%

 Finally, we look at the vertices with an even distance to $\mathcal G^{(2)}_{n,\mu}$. Recall that the number of vertices at distance two or at distance greater than \smash{$K_\gamma := 2 \lfloor\frac{1-\gamma}{2\gamma-1}\rfloor+2$} is of order $o(sn)$. %Hence we only need to consider distances four to $K_\gamma$. If we c
 Choose \smash{$\tilde s = \mu  \beta(n)^{_{\frac{2}{2\gamma-1}}}$} so that $\lfloor\tilde s n \rfloor= m$. Using 
 %the upper bound 
 (\ref{15}) we get
\begin{align*}
    \mathbb E \sum_{k=2}^\infty \#\big(\mathcal{N}_{2k}(\mathcal G_{n,\mu}^{(2)})\setminus \mathcal{N}_{2k+1}(\mathcal G_{n,\mu}^{(2)})\big)\le   \sum_{l=2}^{K_\gamma} \mathbb{E}[X_{2l}] + o(sn) \le  \sum_{l=2}^{K_\gamma} \sum_{0<x_{2l}\le m} \phi_{2l} x_{2l}^{-\gamma} + o(sn)
    \\ \le  c m^{1-\gamma}\beta(n)^2 \log\left(\tfrac{1}{\beta(n)}\right) n ^\gamma + o(sn)= c n \beta(n)^{\tfrac{2\gamma}{2\gamma-1}}\log\left(\tfrac{1}{\beta(n)}\right)+o(sn)\,,
\end{align*}
which is $o(sn)$. 
The result now follows from Markov's inequality.
\end{proof}

\end{document}
